\chapter{Introducción}
Las lesiones cutáneas son alteraciones visibles de la piel que pueden originarse por factores genéticos, ambientales o como manifestación de enfermedades sistémicas. Su clasificación entre benignas y malignas es clínicamente fundamental, ya que determina el curso del tratamiento y el pronóstico del paciente. A pesar de su alta prevalencia a nivel mundial, su diagnóstico oportuno sigue siendo un reto, especialmente en regiones con acceso limitado a servicios especializados. Este proyecto propone el desarrollo de un prototipo de apoyo diagnóstico basado en aprendizaje autosupervisado, orientado a fortalecer las capacidades de identificación de lesiones cutáneas y reducir la dependencia de datos etiquetados en contextos clínicos con recursos limitados como Santander.
